\section*{Capítulo II, Problema VI}

\textbf{Problema:}

6. ¿Cuáles de los siguientes pares de rectas son perpendiculares?

\begin{itemize}
    \item[(a)] $3x - 5y = 1$ y $2x + y = 2$
    \item[(b)] $2x + 7y = 1$ y $x - y = 5$
    \item[(c)] $3x - 5y = 1$ y $5x + 3y = 7$
    \item[(d)] $-x + y = 2$ y $x + y = 9$
\end{itemize}

\textbf{Solución:}

Dos rectas del plano son perpendiculares si el producto de sus pendientes es $-1$.

\begin{itemize}
    \item[(a)]
    \[
        3x - 5y = 1 \implies y = \frac{3}{5}x - \frac{1}{5},
        \qquad
        2x + y = 2 \implies y = -2x + 2.
    \]
    Las pendientes son $m_1 = \frac{3}{5}$ y $m_2 = -2$. Como
    \[
        m_1m_2 = \frac{3}{5}(-2) = -\frac{6}{5} \neq -1,
    \]
    estas rectas no son perpendiculares.

    \item[(b)]
    \[
        2x + 7y = 1 \implies y = -\frac{2}{7}x + \frac{1}{7},
        \qquad
        x - y = 5 \implies y = x - 5.
    \]
    Las pendientes son $m_1 = -\frac{2}{7}$ y $m_2 = 1$. Entonces
    \[
        m_1m_2 = -\frac{2}{7} \neq -1,
    \]
    así que tampoco son perpendiculares.

    \item[(c)]
    \[
        3x - 5y = 1 \implies y = \frac{3}{5}x - \frac{1}{5},
        \qquad
        5x + 3y = 7 \implies y = -\frac{5}{3}x + \frac{7}{3}.
    \]
    Las pendientes son $m_1 = \frac{3}{5}$ y $m_2 = -\frac{5}{3}$. Por tanto,
    \[
        m_1m_2 = \frac{3}{5}\left(-\frac{5}{3}\right) = -1.
    \]
    Este par sí es perpendicular.

    \item[(d)]
    \[
        -x + y = 2 \implies y = x + 2,
        \qquad
        x + y = 9 \implies y = -x + 9.
    \]
    Las pendientes son $m_1 = 1$ y $m_2 = -1$. Entonces
    \[
        m_1m_2 = -1.
    \]
    Este par también es perpendicular.
\end{itemize}

En conclusión, los pares perpendiculares son \textbf{(c)} y \textbf{(d)}.